\section{Discussion}\label{sec:discussion}

\paragraph{What the within--between gap means.}
The headline result---a strong positive between-country association
between mean GenAI exposure and institutional trust, paired with a
within-country effect indistinguishable from zero---is the canonical
\citet{SchmidtCatranFairbrother2016} configuration in which a naive
specification that constrains $\gamma_W = \gamma_B$ would
mis-estimate both. Read substantively, it says that the
cross-sectional positive correlation reflects everything that makes
a country a high-AI-exposure country (income per capita, share of
the workforce in services and knowledge occupations, density of
high-skill institutions) rather than the AI exposure itself. The
correlation flows from the third variable that drives both. Within a
given country over 2012--2024, the variation in country-mean exposure
that is identifiable separately from a global time trend is too
small, or too uncorrelated with trust, to move the outcome.

\paragraph{Why this isn't the robot-exposure story.}
The single-country robot-exposure literature
\citep{ColantoneStanig2018,FreyBergerChen2018,AnelliColantoneStanig2021}
identifies effects from local-labour-market shocks within a single
national institutional environment, where between-region variation in
exposure is plausibly orthogonal to baseline trust. Pooling
cross-nationally exposes that the same exercise across the European
panel cannot rule out the third-variable confound. The literature's
mechanism (structural grievance, performance disappointment) may
still hold at the local level within countries; this paper does not
adjudicate that. What it shows is that when you scale the analysis
to the multi-country panel that is the natural unit for European
political-attitude research, the within-between distinction reshapes
the answer.

\paragraph{Pre-registered null on H4.}
The H4 cross-level interaction was deliberately specified with
both substitution and curator readings, plus a precise null. The
data return the null. The within-country exposure effect (such as
it is) is not heterogeneous in the individual education gradient,
which falsifies both substitution and curator framings on the
within-country margin. This is a strong null---the interaction is
small in absolute value and tightly estimated, not merely
underpowered.

\paragraph{Endogeneity.}
The within-between specification handles time-invariant
country-level confounders---welfare-state arrangements, baseline
political culture, geography---but not time-varying ones.
A confounder that moves with country-year exposure and also moves
trust in the same direction (e.g.\ contemporaneous shifts in media
environment, immigration policy, or party-system fragmentation)
would not be absorbed by the country random intercept. Reverse
causation is also plausible: declining trust may accelerate AI
adoption via labour-market deregulation or reduced public investment
in retraining, rather than the other way around. The
cross-sectional estimate $\hat\gamma_B$ in particular should be
read as an association \emph{within a multilevel structure}, not a
causal effect.

\paragraph{Limitations.}
Three are worth flagging beyond the endogeneity caveat above.
\emph{First}, the L2 country-year VPC of $0.022$ in M0 is at the
soft floor of where a within--between decomposition is well-powered
to distinguish $\gamma_W$ from $\gamma_B$. The Wald rejection of
equality is driven primarily by the precision of $\hat\gamma_W$
(near-zero point estimate with reasonable standard error) rather
than by a strong identifying signal in $\hat\gamma_W$.
\emph{Second}, the leave-one-out attenuation of $\hat\gamma_B$ from
$+11.24$ to $+5.44$ implies that the headline between-effect is
inflated by the mechanical correlation between an individual's own
exposure and the country-year average. The leave-one-out estimate
is the conceptually cleaner one; future work should use it as the
primary specification.
\emph{Third}, three robustness avenues are deferred:
the alternative occupation-exposure measures (Webb 2020, Felten
2021, Felten 2023, Eloundou 2024), which require an
ISCO-08$\leftrightarrow$SOC-2010 crosswalk and additional public
downloads; the Kenward-Roger correction on L3 standard errors via
\texttt{pymer4}/\texttt{lmerTest}; and the multiple imputation of
\texttt{hinctnta} missingness.

\paragraph{Conclusion.}
The diffusion of generative AI capabilities since 2022 has not, in
the period this panel covers, moved European institutional trust
within countries in the way single-country robot-exposure studies
might lead one to expect. The cross-sectional positive association
between AI exposure and trust does not survive the
within-between decomposition. Whether this changes as the GenAI
capability frontier diffuses further into white-collar labour
markets is an open empirical question that the next ESS round
(R12, fielding 2026) will let us begin to answer.
